% Options for packages loaded elsewhere
\PassOptionsToPackage{unicode}{hyperref}
\PassOptionsToPackage{hyphens}{url}
\PassOptionsToPackage{dvipsnames,svgnames,x11names}{xcolor}
%
\documentclass[
  10pt,
  a4paper,
]{article}
\usepackage{amsmath,amssymb}
\usepackage{setspace}
\usepackage{iftex}
\ifPDFTeX
  \usepackage[T1]{fontenc}
  \usepackage[utf8]{inputenc}
  \usepackage{textcomp} % provide euro and other symbols
\else % if luatex or xetex
  \usepackage{unicode-math} % this also loads fontspec
  \defaultfontfeatures{Scale=MatchLowercase}
  \defaultfontfeatures[\rmfamily]{Ligatures=TeX,Scale=1}
\fi
\usepackage{lmodern}
\ifPDFTeX\else
  % xetex/luatex font selection
  \setmainfont[]{Liberation Serif}
  \setsansfont[]{Liberation Sans}
  \setmonofont[]{DejaVu Sans Mono}
\fi
% Use upquote if available, for straight quotes in verbatim environments
\IfFileExists{upquote.sty}{\usepackage{upquote}}{}
\IfFileExists{microtype.sty}{% use microtype if available
  \usepackage[]{microtype}
  \UseMicrotypeSet[protrusion]{basicmath} % disable protrusion for tt fonts
}{}
\makeatletter
\@ifundefined{KOMAClassName}{% if non-KOMA class
  \IfFileExists{parskip.sty}{%
    \usepackage{parskip}
  }{% else
    \setlength{\parindent}{0pt}
    \setlength{\parskip}{6pt plus 2pt minus 1pt}}
}{% if KOMA class
  \KOMAoptions{parskip=half}}
\makeatother
\usepackage{xcolor}
\usepackage[margin=22mm]{geometry}
\setlength{\emergencystretch}{3em} % prevent overfull lines
\providecommand{\tightlist}{%
  \setlength{\itemsep}{0pt}\setlength{\parskip}{0pt}}
\setcounter{secnumdepth}{-\maxdimen} % remove section numbering
\ifLuaTeX
\usepackage[bidi=basic]{babel}
\else
\usepackage[bidi=default]{babel}
\fi
\babelprovide[main,import]{korean}
\ifPDFTeX
\else
\babelfont{rm}[]{Liberation Serif}
\fi
% get rid of language-specific shorthands (see #6817):
\let\LanguageShortHands\languageshorthands
\def\languageshorthands#1{}
\usepackage{xeCJK}
\setCJKmainfont{Noto Serif CJK KR}
\setCJKsansfont{Noto Sans CJK KR}
\setCJKmonofont{Noto Sans Mono CJK KR}
\xeCJKsetup{CJKspace=true}
\usepackage{xurl}
\usepackage{fvextra}
\usepackage{fancyhdr}
\AtBeginDocument{\DefineVerbatimEnvironment{verbatim}{Verbatim}{breaklines=true,breakanywhere=true,fontsize=\small}}
\setlength{\emergencystretch}{4em}
\setlength{\headheight}{15pt}
\pagestyle{fancy}\fancyhf{}
\fancyhead[L]{\small 중성자별 구조와 비방사 진동}
\fancyfoot[C]{\thepage}
\renewcommand{\headrulewidth}{0.3pt}
\setcounter{tocdepth}{2}
\ifLuaTeX
  \usepackage{selnolig}  % disable illegal ligatures
\fi
\usepackage{bookmark}
\IfFileExists{xurl.sty}{\usepackage{xurl}}{} % add URL line breaks if available
\urlstyle{same}
\hypersetup{
  pdftitle={A03 · Euler와 RK4: 단계와 전역 오차},
  pdflang={ko-KR},
  colorlinks=true,
  linkcolor={Maroon},
  filecolor={Maroon},
  citecolor={Blue},
  urlcolor={Blue},
  pdfcreator={LaTeX via pandoc}}

\title{A03 · Euler와 RK4: 단계와 전역 오차}
\usepackage{etoolbox}
\makeatletter
\providecommand{\subtitle}[1]{% add subtitle to \maketitle
  \apptocmd{\@title}{\par {\large #1 \par}}{}{}
}
\makeatother
\subtitle{과제 큐시트}
\author{}
\date{}

\begin{document}
\maketitle

\setstretch{1.12}
\textbf{선행:} A02　 \textbf{독립 작업:} 3--5시간 (계획치)

\subsection{적분값 하나가 맞았는데도 왜 다른 예제를
푸는가?}\label{uxc801uxbd84uxac12-uxd558uxb098uxac00-uxb9deuxc558uxb294uxb370uxb3c4-uxc65c-uxb2e4uxb978-uxc608uxc81cuxb97c-uxd478uxb294uxac00}

우변이 상태에만 의존하는 예제에서는 중간 좌표를 잘못 전달해도 오류가
드러나지 않을 수 있다. 좌표 자체가 우변에 들어가는 예제가 필요한 이유를
먼저 적자.

먼저 한두 문장으로 답한 뒤 아래 설명을 읽자. 첫 예측이 틀렸거나 아직
모르겠다고 적는 것은 감점하지 않는다.

\subsection{필요한 관계와 읽을
위치}\label{uxd544uxc694uxd55c-uxad00uxacc4uxc640-uxc77duxc744-uxc704uxce58}

\(y_{n+1}=y_n+h(k_1+2k_2+2k_3+k_4)/6\). 매끄럽고 안정한 유한 구간에서
Euler와 RK4 전역 오차는 각각 \(O(h)\), \(O(h^4)\)이다.

\begin{itemize}
\tightlist
\item
  \path|notes/chapters/ch02_numerics.md|
\item
  \path|learner/src/numerics/ns_rk4.f90|
\item
  \path|tools/numerics_demo.py|
\end{itemize}

\subsection{계산으로
옮기기}\label{uxacc4uxc0b0uxc73cuxb85c-uxc62euxae30uxae30}

기본 경로는 제공 함수의 호출과 해석이다. 아래 \texttt{learner} 시험은
같은 물리 함수를 직접 작성한 경우에 실행한다. 제공 예제를 검사할 때에는
명시적으로 \texttt{calculation}을 선택한다.

\begin{enumerate}
\def\labelenumi{\arabic{enumi}.}
\tightlist
\item
  해석해를 기준으로 Euler·RK4의 한 단계를 손으로 확인한다.
\item
  제공 방법 또는 자신이 작성한 함수를 사용하여 N=16,32,64에서 오차를
  비교한다.
\item
  비자율 문제와 조화진동자에서 좌표 전달·진폭·위상 오차를 분리한다.
\end{enumerate}

\newpage

\subsection{실행해서 무엇을 볼
것인가}\label{uxc2e4uxd589uxd574uxc11c-uxbb34uxc5c7uxc744-uxbcfc-uxac83uxc778uxac00}

패키지 루트에서 실행한다. \texttt{-\/-plan}은 경로만 확인하며 계산
성공을 뜻하지 않는다.

\begin{verbatim}
python tools/run_assignment.py A03 --plan
python tools/run_assignment.py A03 --run \
  --workspace learner --output ../A03_result.json
\end{verbatim}

\subsection{결과를
설명하기}\label{uxacb0uxacfcuxb97c-uxc124uxba85uxd558uxae30}

오차표에서는 단계 수를 바꾼 비율과 오차가 더 이상 줄지 않는 구간을
구별한다. 조화진동자는 진폭뿐 아니라 위상도 비교해야 한다.

\textbf{남길 증거.} 단계표·수렴 표 또는 그림; 비자율 방정식 검사; RK4가
에너지를 정확히 보존하지 않는 이유

\textbf{판단 기준.} 2점 예: 세 해상도 값과 오차비를 제시하며 이 비를
엄밀한 전역오차 상한으로 부르지 않는다. 최초 예측의 정오보다, 결과를 본
뒤 이유를 고치고 다른 입력에도 적용하는지를 본다.

\subsection{막혔을 때}\label{uxb9c9uxd614uxc744-uxb54c}

\textbf{힌트 1.} k에는 h가 포함되는지 API 정의를 확인한다.

\textbf{힌트 2.} y′=x³를 사용하면 단계의 x 인수를 값 계산과 독립적으로
확인할 수 있다.

\subsection{다음 만남에 다시
묻기}\label{uxb2e4uxc74c-uxb9ccuxb0a8uxc5d0-uxb2e4uxc2dc-uxbb3buxae30}

다음 회차에는 처음 보지 않은 적분 구간을 정하고 어떤 해석해를 기준으로
삼을지 스스로 고른다. 차수 이름만 답하지 않는다. 먼저 짧게 답한 뒤
본문과 자신의 기록을 열어 확인한다.

\textbf{범위.} 차수는 안정성이나 임의 step 크기의 정확도를 보증하지
않는다.

\end{document}
